\documentclass[11pt]{amsart}
\usepackage{amsmath,amssymb,amsthm}
\usepackage{url}
\emergencystretch=3em

\newtheorem{theorem}{Theorem}[section]
\newtheorem{proposition}[theorem]{Proposition}
\newtheorem{corollary}[theorem]{Corollary}
\newtheorem{lemma}[theorem]{Lemma}
\theoremstyle{remark}
\newtheorem{remark}[theorem]{Remark}
\newtheorem{question}[theorem]{Question}

\newcommand{\Aa}{A_{\alpha}}
\newcommand{\dd}{\boldsymbol{d}}
\newcommand{\one}{\boldsymbol{1}}
\newcommand{\E}{\mathcal{E}}
\newcommand{\QQ}{\mathbb{Q}}

\begin{document}
\title[Equality in average-neighbour-degree bounds]{Equality in Nikiforov's average-neighbour-degree bounds for $A_\alpha$}
\author{Aleksei Kudriashov}
\address{Athanor lab, Bali, Indonesia}
\thanks{ORCID: \url{https://orcid.org/0009-0006-8885-5338}.}
\date{15 November 2026 (draft)}
\subjclass[2020]{05C50, 15A18, 05C75}
\keywords{$A_\alpha$-matrix, spectral radius, average neighbour degree, equality case, degree sequence, computational enumeration}

\begin{abstract}
For a graph without isolated vertices, Nikiforov bounded the spectral radius
of $A_\alpha=\alpha D+(1-\alpha)A$ by the minimum and maximum of
$\alpha d(v)+(1-\alpha)m(v)$, where $m(v)$ is the average degree of the
neighbours of $v$. We study equality for connected, nonregular graphs and
$0<\alpha<1$. Such a graph has at most one equality parameter; it is
rational, belongs to $(0,1/2]$, and forces $A_\alpha$ to be singular.
At $\alpha=1/2$ the graphs are precisely the nonregular semiregular
bipartite graphs. We give an integral recognition criterion and a complete
classification when exactly two degrees occur. Conversely, every rational
number in $(0,1/2]$ occurs. For every $k\geq 2$ an explicit connected
$2k$-vertex graph with $k$ distinct degrees attains equality at $\alpha=1/3$;
thus no absolute bound on degree variety is possible. An exact
isomorphism-class census through ten vertices, independently generated by
nauty and checked with rational arithmetic, yields $599$ irregular cases
at order ten. We distinguish proofs from finite computer certificates,
and delimit the prior-art search rather than asserting publication novelty.
\end{abstract}
\maketitle

\section{Introduction and the bounds}
All graphs in the theoretical statements below are finite, simple, undirected,
connected and of order $n\geq2$, unless specified otherwise. Let $A$ and $D$
denote the adjacency and degree matrices, and put
\[
 A_\alpha(G)=\alpha D+(1-\alpha)A,\qquad 0<\alpha<1.
\]
Write $d(v)$ for the degree, $s(v)=\sum_{w\sim v}d(w)$ for the
neighbour-degree sum, $m(v)=s(v)/d(v)$ for the average neighbour degree,
$\dd=(d(v))_{v\in V(G)}$, and $\rho_\alpha=\rho(A_\alpha(G))$.
Connectedness ensures $d(v)\geq1$.

Nikiforov~\cite[Proposition~20, (20)--(21)]{Nikiforov2017} established,
for graphs without isolated vertices,
\begin{align}
 \rho_\alpha&\leq\max_{v\in V(G)}
       \left\{\alpha d(v)+(1-\alpha)m(v)\right\},\tag{20}\label{eq:upper}\\
 \rho_\alpha&\geq\min_{v\in V(G)}
       \left\{\alpha d(v)+(1-\alpha)m(v)\right\}.\tag{21}\label{eq:lower}
\end{align}
His proposition also proves that, for connected graphs and
$1/2<\alpha<1$, equality occurs precisely for regular graphs. Just before
Problem~22, Nikiforov suggested that apart from $\alpha=0$ and $1/2$,
regularity might be necessary. We examine this informal expectation, not the
valid inequalities or the separate edge bounds (22)--(23) in his
Corollary~21. For example, $K_1\vee K_{2,2}$ attains both bounds at
$\alpha=1/4$ (Section~\ref{sec:examples}).

The lineage matters. Merris's degree-plus-average-neighbour-degree bound
for the Laplacian~\cite{Merris1998} and Feng and Yu's signless-Laplacian
equality theorem~\cite[Lemma~2.3]{FengYu2009} precede Nikiforov's
$A_\alpha$ formulation. The characterization at $\alpha=1/2$
is due to Feng and Yu. In later work,
Liu, Das and Shu~\cite[Lemmas~2.9--2.10 and Theorems~5.4--5.6]{LiuDasShu2020}
consider related degree conditions and equality statements; in particular,
Theorem~5.4 concerns the \emph{edge} bounds, and Theorem~5.6 addresses
components in the $\alpha>1/2$ range, rather than providing the
classification developed below. Wang, Wong and
Tian~\cite{WangWongTian2020} give a different edge-degree spectral bound
and associated regular/semiregular cases. The adjacency endpoint $\alpha=0$
is related to the established theory of harmonic graphs, for which
$A\dd=\lambda\dd$; see, for example,
Petrovi\'c, Borovi\'canin and Radosavljevi\'c~\cite{Petrovic2006}.
Ghebleh, Al-Yakoob, Kanso and Stevanovi\'c~\cite{GheblehEtAl2025}
construct graphs with arbitrarily many distinct degrees and precisely two
main \emph{adjacency} eigenvalues; in particular their harmonic case
satisfies $A\dd=\lambda\dd$ at $\alpha=0$. This is a different
condition from $A_\alpha\dd=\rho_\alpha\dd$ at a fixed positive
$\alpha$. The publication search is bounded; see
Section~\ref{sec:priority}.

\section{The equality equation and structural restrictions}
\begin{lemma}[Perron row-sum criterion]\label{lem:rows}
For a connected graph $G$, equality in either \eqref{eq:upper} or
\eqref{eq:lower} holds if and only if equality holds in both, if and only if
there is a scalar $c$ such that
\begin{equation}\label{eq:eigenvector}
 A_\alpha\dd=c\dd.
\end{equation}
In that event $c=\rho_\alpha$ and
\begin{equation}\label{eq:local}
 \alpha d(v)^2+(1-\alpha)s(v)=c\,d(v)
 \qquad(v\in V(G)).
\end{equation}
\end{lemma}
\begin{proof}
The matrix $B=D^{-1}A_\alpha D$ is nonnegative, irreducible and similar to
$A_\alpha$. Its $v$th row sum equals
$\alpha d(v)+(1-\alpha)s(v)/d(v)$. Let $y$ be a strictly positive left
Perron vector of $B$, so $y^TB=\rho_\alpha y^T$. If the spectral radius
is either the minimum or maximum row sum, then the identity
\[
 \rho_\alpha\sum_v y_v
  =y^TB\one=\sum_v y_v\left(\alpha d(v)
                  +(1-\alpha)\frac{s(v)}{d(v)}\right)
\]
forces \emph{every} row sum to equal $\rho_\alpha$, because every $y_v>0$.
Conversely, constant row sums give $B\one=c\one$; equivalently
$A_\alpha\dd=c\dd$. Since $\dd>0$, Perron--Frobenius implies
$c=\rho_\alpha$, making both bounds equalities. Multiplying the $v$th
coordinate by $d(v)$ gives \eqref{eq:local}.
\end{proof}

\begin{theorem}[Structure and exact recognition]\label{thm:structure}
Let $G$ be connected of order $n\geq2$.
\begin{enumerate}
\item If $G$ is regular, equality in \eqref{eq:upper} and
\eqref{eq:lower} holds for every $0<\alpha<1$.
\item If $G$ is nonregular, there is at most one equality parameter;
if it exists, it is rational and lies in $(0,1/2]$. At $\alpha=1/2$,
the nonregular equality graphs are exactly the connected, semiregular
bipartite graphs (degree constant on each bipartition class).
\item More precisely, fix a maximum-degree vertex $x$ and a
minimum-degree vertex $y$, and set
\[
 a=d(x)=\Delta,\quad b=d(y)=\delta,\quad
 H=(a-b)ab,\quad K=a s(y)-b s(x).
\]
Then $G$ is a nonregular equality graph for some $0<\alpha<1$
if and only if
\begin{equation}\label{eq:integer-criterion}
 0<K\leq H,\qquad
 a\bigl(Kd(v)^2+Hs(v)\bigr)
  =d(v)\bigl(Ka^2+Hs(x)\bigr)
 \quad\text{for every }v.
\end{equation}
The unique parameter is
\begin{equation}\label{eq:parameter}
 \alpha=\frac{K}{H+K}.
\end{equation}
In lowest terms its denominator is at most
$2(n-2)(n-1)^2$.
\item At any such nonregular equality parameter, $A_\alpha$ is
singular, and $D^{-1/2}AD^{-1/2}$ has eigenvalue
$-\alpha/(1-\alpha)$.
\end{enumerate}
\end{theorem}
\begin{proof}
For a $r$-regular graph, $\dd=r\one$ and
$A_\alpha\dd=r\dd$; Lemma~\ref{lem:rows} proves (1).
Suppose henceforth that $G$ is nonregular and equality holds.
By Lemma~\ref{lem:rows}, $A_\alpha\dd=c\dd$. But
$A_\alpha\one=\dd$, whence
\begin{equation}\label{eq:kernel}
 A_\alpha(\dd-c\one)=0.
\end{equation}
The vector $\dd-c\one$ is nonzero, since otherwise $G$ would be
regular. Thus $A_\alpha$ is singular.

Set $Q=D+A$. The identity
\[
 A_\alpha=(1-\alpha)Q+(2\alpha-1)D,
 \qquad z^TQz=\sum_{uv\in E(G)}(z_u+z_v)^2
\]
shows that $A_\alpha$ is positive definite when $\alpha>1/2$:
$D$ is positive definite and both coefficients are nonnegative,
with $2\alpha-1>0$. This contradicts singularity. Hence
$\alpha\leq1/2$.

If $\alpha=1/2$, equation \eqref{eq:kernel} says that
$Q(\dd-c\one)=0$. For any $z$, the displayed sum of squares
vanishes precisely when $z_u=-z_v$ on each edge. In a connected
nonbipartite graph an odd cycle forces $z=0$, and connectedness then
forces zero everywhere. Thus $Q$ is nonsingular in that case.
For a connected bipartite graph with parts $U,W$, its nullspace is
one-dimensional, spanned by the vector $\varepsilon$ taking values
$1$ on $U$ and $-1$ on $W$. Consequently
$\dd-c\one=t\varepsilon$ and the degrees are constant, say $r$
on $U$ and $s$ on $W$. Conversely, for such a graph,
$Q\dd$ has entry $r^2+rs=(r+s)r$ on $U$ and
$s^2+sr=(r+s)s$ on $W$, so
$A_{1/2}\dd=(r+s)\dd/2$. Lemma~\ref{lem:rows}
proves equality. The nonregular condition is $r\ne s$.

For the remaining assertions, compare \eqref{eq:local} at the
vertices $x$ and $y$ after dividing by their degrees:
\[
 \alpha(a-b)=(1-\alpha)
       \left(\frac{s(y)}b-\frac{s(x)}a\right).
\]
Multiplication by $ab$ gives $\alpha H=(1-\alpha)K$.
Here $H>0$, so $K>0$ and $\alpha=K/(H+K)$.
The already established inequality $\alpha\leq1/2$
gives $K\leq H$. This formula proves both uniqueness and rationality.
Substitution into the row sum gives, for every vertex,
\[
 \alpha d(v)+(1-\alpha)\frac{s(v)}{d(v)}
   =\frac{Kd(v)^2+Hs(v)}{(H+K)d(v)}.
\]
All row sums coincide exactly when the integer equations in
\eqref{eq:integer-criterion} hold; the displayed formula also
proves the converse by Lemma~\ref{lem:rows}. Since
$a,b\leq n-1$ and $a-b\leq n-2$,
$H\leq(n-2)(n-1)^2$; with $K\leq H$, the unreduced denominator
$H+K\leq2(n-2)(n-1)^2$, which bounds the reduced one.
Finally,
\[
 A_\alpha=D^{1/2}
 \left(\alpha I+(1-\alpha)D^{-1/2}AD^{-1/2}\right)D^{1/2}.
\]
Its singularity implies the asserted normalized-adjacency eigenvalue,
completing all four parts.
\end{proof}

\section{Exactly two degrees and regular joins}\label{sec:twodegrees}
\begin{theorem}[Two-degree classification]\label{thm:two}
Suppose $G$ has precisely two distinct degrees $a>b$.
Equality at some $0<\alpha<1$ forces the partition into the two degree
classes to be equitable: each degree-$a$ vertex has the same number
$x$ of degree-$a$ neighbours and each degree-$b$ vertex has the same
number $y$ of degree-$a$ neighbours. Conversely, given a connected
two-degree graph with this equitable partition, put
\[
 h=\frac yb-\frac xa.
\]
There is an equality parameter in $(0,1)$ if and only if $h>0$;
in that event it is unique and equals
\begin{equation}\label{eq:h}
 \alpha=\frac{h}{1+h}.
\end{equation}
Automatically $0<h\leq1$, and $h=1$ if and only if the graph is
semiregular bipartite with the two degree classes as its parts.
\end{theorem}
\begin{proof}
At a vertex of degree $d\in\{a,b\}$ having $t$ neighbours of
degree $a$, its neighbour-degree sum is $ta+(d-t)b$ and hence
\[
 m=b+(a-b)\frac{t}{d}.
\]
If equality holds, all row sums are constant by Lemma~\ref{lem:rows}.
For a fixed $d$ the term $\alpha d$ is fixed; because
$1-\alpha>0$, $m$ must be fixed. As $a-b>0$, $t$ is fixed on
each degree class, giving the stated $x,y$.
The two row sums therefore differ (high minus low) by
\[
 (a-b)\left[\alpha+(1-\alpha)
                  \left(\frac xa-\frac yb\right)\right].
\]
They coincide precisely when $\alpha=(1-\alpha)h$, that is,
\eqref{eq:h}. This lies in $(0,1)$ exactly when $h>0$;
then Lemma~\ref{lem:rows} proves equality and uniqueness.
Since $0\leq x/a\leq1$ and $0\leq y/b\leq1$, we have $h\leq1$.
Equality $h=1$ requires $y=b$ and $x=0$, and conversely these
values give $h=1$. Thus every low-degree vertex has all its
neighbours in the high-degree class, and no high-degree vertex
has a neighbour in its own class. Both degree classes are
independent, so $G$ is semiregular bipartite. The reverse implication
follows from its cross-class adjacency counts.
\end{proof}

\begin{proposition}[Uniform independent blow-ups]\label{prop:blow}
If a connected nonregular graph $G$ attains equality at $\alpha$,
then for every integer $t\geq1$ the graph obtained by replacing
any vertex by $t$ independent clones and any edge by all $t^2$
possible clone-to-clone edges is connected, attains equality at the
same parameter, has the same number of distinct degrees, and has
spectral radius $t\rho_\alpha(G)$.
\end{proposition}
\begin{proof}
For a clone of $v$, its degree is $td(v)$. For every original
neighbour $w$ there are $t$ clones of $w$, each of degree $td(w)$,
so its neighbour-degree sum is $t^2s(v)$ and its average neighbour
degree is $tm(v)$. Therefore every new row sum is
$t[\alpha d(v)+(1-\alpha)m(v)]=t\rho_\alpha(G)$.
The blow-up is connected because any path in $G$ lifts between
clones of distinct vertices; clones of one vertex are linked by a
length-two path through any neighbour. It remains nonregular, and
multiplying all degrees by $t$ preserves their number of distinct
values. Lemma~\ref{lem:rows} completes the proof.
\end{proof}

\begin{theorem}[Complete joins of regular pieces]\label{thm:joins}
Suppose $G$ is the complete join of finitely many internally regular
graphs (possibly of different orders and internal degrees).
If $G$ attains equality at some $0<\alpha<1$, then it has at most
two distinct vertex degrees.
\end{theorem}
\begin{proof}
Let $n=|V(G)|$ and $S=\sum_{v\in V(G)}d(v)$.
For a vertex $v$ of total degree $d$ lying in a regular join piece,
all vertices not adjacent to $v$ (including $v$ itself) belong
to that piece and have total degree $d$: indeed the internal
degree is constant within the piece and every vertex there sees
all vertices outside the piece. There are exactly $n-d$ such
vertices. Subtracting their contribution from $S$ gives
$s(v)=S-(n-d)d$, hence
\[
 m(v)=d-n+\frac Sd.
\]
If the common row sum is $c$, multiplication of
$c=d-(1-\alpha)n+(1-\alpha)S/d$ by $d$ yields
\[
 d^2-\bigl(c+(1-\alpha)n\bigr)d+(1-\alpha)S=0.
\]
Every occurring degree is a real root of this single nonzero
quadratic with leading coefficient one. There are at most two.
\end{proof}

\section{Attainable parameters and many degrees}\label{sec:examples}
First, simple examples already contradict the informal expectation
preceding Nikiforov's Problem~22. For $K_t\vee\overline{K_m}$,
$t,m\geq2$, the two degrees are $a=t+m-1$ and $b=t$, with
$x=t-1$ and $y=t$. Theorem~\ref{thm:two} gives
\begin{equation}\label{eq:split}
 \alpha=\frac{m}{2m+t-1}.
\end{equation}
The graph $K_1\vee K_{2,2}$ is another connected, nonregular example:
the degree-$4$ vertex has neighbour-degree sum $12$ and the four
degree-$3$ vertices each have sum $10$. At $\alpha=1/4$ the two
row sums are both $13/4$, so $\rho_{1/4}=13/4$ by
Lemma~\ref{lem:rows}. A published special case appears in Alhevaz, Baghipur, Ganie and
Das~\cite[Remark~2.4]{AlhevazEtAl2023}: for
$K_4-e=K_2\vee\overline{K_2}$, they give the two expressions
$3-\alpha$ and $7/3+(2/3)\alpha$ in Nikiforov's vertexwise bound.
They coincide at $\alpha=2/5$, implicitly giving equality in both
(20)--(21) by Lemma~\ref{lem:rows}. Their remark compares upper
bounds and does not explicitly frame the example as a counterexample
to Nikiforov's expectation. The general family
$K_t\vee\overline{K_m}$ with parameter \eqref{eq:split}, and that
framing, were not found in two independent bounded audits and the
author's search; this is not a publication-priority claim.

\begin{theorem}[All interior rational parameters]\label{thm:rationals}
Among connected nonregular graphs, the parameters $0<\alpha<1$
at which equality holds in (20)--(21) form precisely
$\QQ\cap(0,1/2]$. The set is dense in $(0,1/2)$, its closure in
$\mathbb R$ is $[0,1/2]$, and its reduced denominators are unbounded.
\end{theorem}
\begin{proof}
Necessity follows from Theorem~\ref{thm:structure}.
For sufficiency, take a rational $p/q$ in lowest terms with
$0<2p<q$, set $r=q-2p\geq1$, and form a graph from a clique
$K_{r+1}$ by attaching $p$ pendant leaves to each clique vertex.
It is connected and nonregular. Its clique vertices have degree
$a=r+p$, with exactly $x=r$ high-degree neighbours; its leaves
have degree $b=1$, with $y=1$ high-degree neighbour.
Theorem~\ref{thm:two} applies, since
\[
 h=\frac11-\frac{r}{r+p}=\frac{p}{r+p}>0,
 \qquad \alpha=\frac{p}{r+2p}=\frac pq.
\]
This uses a direct algebraic criterion, not a numerical search.
For $p/q=1/2$, the star $K_{1,2}$ is semiregular bipartite
and nonregular, so Theorem~\ref{thm:structure} gives equality.
Every rational in the stated interval is therefore realized.
Density and the closure assertion are the elementary density of
rational numbers, and taking, for example, $\alpha=1/q$ for
$q\geq3$ proves that the reduced denominators are unbounded.
\end{proof}

\begin{theorem}[The staircase family]\label{thm:stairs}
For every integer $k\geq2$ there is a connected simple graph
$T_k$ on $2k$ vertices with exactly $k$ distinct degrees
$1,2,\ldots,k$ that attains equality in (20)--(21) at
$\alpha=1/3$. Its spectral radius is $(2k+1)/3$.
Consequently the number of distinct degrees in connected
nonregular equality graphs is unbounded.
\end{theorem}
\begin{proof}
Take vertices $(i,b)$, $1\leq i\leq k$, $b\in\{0,1\}$.
For distinct levels $i,j$ join $(i,b)$ to $(j,b)$ if and only if
$i+j\geq k+1$; the two endpoints must be in the same layer.
In addition, insert the rung $(i,0)(i,1)$ if and only if
$2i\geq k+1$. There are no other edges. Since each pair of
levels contributes at most one same-layer edge per layer,
and each rung has distinct endpoints, this is a simple graph.

For fixed $i$, the eligible levels are exactly
$j=k+1-i,k+2-i,\ldots,k$, of which there are $i$.
If $j\ne i$, the unique neighbour at level $j$ is $(j,b)$;
if $j=i$, eligibility is $2i\geq k+1$ and the unique
neighbour at that level is $(i,1-b)$ across the rung.
Thus each vertex at level $i$ has degree exactly $i$.
In particular, there are precisely $k$ distinct degrees.
Every vertex $(i,b)$ with $i<k$ is adjacent to $(k,b)$,
and $(k,0)$ is joined to $(k,1)$ by a rung, so $T_k$
is connected.

The sum of the degrees of the neighbours of either vertex
at level $i$ is
\[
 s_i=\sum_{j=k+1-i}^{k}j
     =\frac{i(2k+1-i)}2.
\]
Accordingly, at $\alpha=1/3$,
\[
 3(A_{1/3}\dd)_{(i,b)}
   =i^2+2s_i=i^2+i(2k+1-i)=(2k+1)i.
\]
Hence $A_{1/3}\dd=((2k+1)/3)\dd$.
The degree vector is positive and the graph is connected,
so Lemma~\ref{lem:rows} identifies this eigenvalue as the
spectral radius and proves equality in both bounds.
As $k$ is arbitrary, the degree-variety conclusion follows.
\end{proof}

The degree partition in Theorem~\ref{thm:stairs} is equitable,
with quotient entries $b_{ij}=1$ if $i+j\geq k+1$ and $0$
otherwise, but an arbitrary equality graph need \emph{not}
have an equitable degree partition. The following explicit
example also documents why the two-degree classification
cannot be extended merely by requiring equitability.

\begin{proposition}[Non-equitable degree partition]\label{prop:noneq}
There exists a connected nine-vertex graph with three distinct
degrees, having equality at $\alpha=1/3$, whose partition into
degree classes is not equitable.
\end{proposition}
\begin{proof}
On vertices $0,\ldots,8$, take the edge set
\[
\begin{split}
\{&(0,3),(1,4),(2,5),(0,6),(1,6),(2,6),
(0,7),(1,7),\\
 &(2,7),(6,7),(0,8),(1,8),(2,8),(3,8),
 (4,8),(5,8)\}.
\end{split}
\]
The degrees, in vertex order, are
$(4,4,4,2,2,2,4,4,6)$, and the neighbour-degree sums are
$(16,16,16,10,10,10,16,16,18)$.
Direct substitution gives
$[d(v)^2+2s(v)]/[3d(v)]=4$ at every vertex.
The graph is connected, since vertex $8$ meets
$0,1,2,3,4,5$ and vertices $6,7$ meet $0,1,2$.
Lemma~\ref{lem:rows} yields $\rho_{1/3}=4$.
Yet degree-$4$ vertex $0$ has respectively $1,2,1$
neighbours of degrees $2,4,6$, while degree-$4$
vertex $6$ has respectively $0,4,0$ such neighbours.
Thus the degree partition is not equitable.
\end{proof}

\section{Exact finite census through ten vertices}\label{sec:census}
We call two graphs equivalent if they are isomorphic. Let $C_n$
denote the number of connected classes, $R_n$ the number of regular
connected classes, and $E_n$ the number of connected nonregular
classes for which equality holds for at least one $\alpha\in(0,1)$.
Theorem~\ref{thm:structure} guarantees that there is exactly
one such $\alpha$ per counted nonregular class.

\begin{theorem}[Computer-assisted census]\label{thm:census}
The exact numbers for $2\leq n\leq10$ are as follows. The columns
$E_n(j)$ count the irregular equality graphs with $j$ distinct
vertex degrees; $E_n(1/2)$ counts those attaining equality at
$\alpha=1/2$.
\begin{center}
\small
\begin{tabular}{r|r r r|r r r r|r}
$n$ & $C_n$ & $R_n$ & $E_n$ & $E_n(2)$ & $E_n(3)$ & $E_n(4)$ & $E_n(5)$ & $E_n(1/2)$\\\hline
2 & 1 & 1 & 0 & 0 & 0 & 0 & 0 & 0\\
3 & 2 & 1 & 1 & 1 & 0 & 0 & 0 & 1\\
4 & 6 & 2 & 3 & 3 & 0 & 0 & 0 & 1\\
5 & 21 & 2 & 7 & 6 & 1 & 0 & 0 & 2\\
6 & 112 & 5 & 19 & 17 & 2 & 0 & 0 & 2\\
7 & 853 & 4 & 23 & 18 & 4 & 1 & 0 & 3\\
8 & 11,117 & 17 & 82 & 77 & 2 & 3 & 0 & 3\\
9 & 261,080 & 22 & 128 & 107 & 19 & 1 & 1 & 5\\
10 & 11,716,571 & 167 & 599 & 538 & 49 & 6 & 6 & 6
\end{tabular}
\end{center}
No irregular equality class with six distinct degrees occurs
at order at most ten. All finite statements here have the
\emph{exact-enumeration, exact-rational-verification} evidence
status described in the proof; they are conditional on the
correctness of the supplied generators, software, and certificate
files, not on floating-point tolerances.
\end{theorem}
\begin{proof}[Computation and exhaustiveness certificate]
We specify both exhaustive enumerations and their equality tests,
so that the finite assertion is reproducible rather than an
unqualified numerical claim. The full source files, output lists,
run instructions and SHA-256 manifests are in
\url{https://github.com/iwasborninbali/paper-aalpha-degree-equality},
under \path{certificates/classify22/},
\path{certificates/guide-checks/}, and
\path{certificates/manydegrees/}; the associated archived research
material is identified by DOI \url{10.5281/zenodo.24239158}.

\emph{Primary independent catalogue, $2\leq n\leq10$.}
The guide generated every connected unlabelled graph with
\texttt{geng} from nauty version 2.9.3~\cite{McKayPiperno2014},
using for each $N=2,\ldots,10$ the command
\begin{center}
\texttt{geng -cq N | python3 eq22.py > eq-nN.log}
\end{center}
from the \texttt{guide-checks} directory. The option
\texttt{-c} restricts to connected graphs and \texttt{-q}
suppresses auxiliary messages. The graphs are serialized in
\texttt{graph6}; for $n\leq10$ its initial character encodes $n$,
and the following bits encode the upper triangle in the order
$(0,1),(0,2),(1,2),(0,3),\ldots,(n-2,n-1)$.
Each \texttt{eq-nN.log} contains one line
\texttt{graph6}, \texttt{alpha}, and \texttt{degree-count} (tab separated) for every
nonregular equality class, followed by the number of
connected, regular and irregular qualifying classes.

For clarity, here is the \emph{entire logical test} of
\texttt{eq22.py}, which uses Python's \texttt{fractions.Fraction},
never floating point. On a graph, compute the integers
$d(v)>0$ and $s(v)$, and the exact fractions
$m(v)=s(v)/d(v)$. Count a constant-degree graph as regular.
Otherwise, reject if two vertices of the same degree have
different $m(v)$. Sort the distinct pairs $(d(v),m(v))$
by degree. Compute the rational slope
$\sigma=(m_2-m_1)/(d_2-d_1)$ from the first two pairs;
reject if $\sigma\geq0$ or if for any further pair
$m_i-m_1\ne\sigma(d_i-d_1)$. Otherwise print
$\alpha=-\sigma/(1-\sigma)$ and the number of distinct
degrees. The test is necessary and sufficient: passing means
$m(v)=c_0+\sigma d(v)$ for every vertex, so
$\alpha d(v)+(1-\alpha)m(v)=(1-\alpha)c_0$
when $\alpha/(1-\alpha)=-\sigma$; conversely constant
row sums force that slope. Lemma~\ref{lem:rows} then
identifies the exact equality cases. The output was
independently audited by \texttt{manydegrees/census.py}, which
decodes every retained \texttt{graph6} string and checks
\eqref{eq:local} with cross-multiplied integers, as well as
connectivity, degree variety, and equitable or non-equitable
neighbour-degree histograms.

\emph{Separate exhaustive catalogue, $2\leq n\leq9$.}
The C11 program \path{classify22/census.c} starts at the
single vertex and recursively adjoins a new vertex along every
nonempty subset of the existing vertices. This covers every
connected graph: every connected graph on at least two
vertices has a vertex whose deletion leaves the graph
connected (choose a leaf of a spanning tree), and induction
applies to the remaining graph. To remove isomorphic duplicates,
the program colors vertices initially by degree, repeatedly
refines the coloring by the counts of neighbours in each color
class until stable, and minimizes the adjacency bit code over
all permutations within the resulting \emph{ordered} invariant
color classes. The stable ordered classes are preserved by
isomorphisms; hence isomorphic inputs have identical minima.
Conversely identical encoded adjacency matrices give an
isomorphism. Its open-addressed hash table retains one
representative per code. For each representative the program
uses the \emph{integer} inequalities and equations
\eqref{eq:integer-criterion}; it never estimates an eigenvalue.
It writes a line
with fields \texttt{W}, \texttt{n}, \texttt{hex-code}, \texttt{alpha}, \texttt{degree-count}
per nonregular case to \path{classify22/census.log},
including all $263$ cases at orders $2$--$9$.
The separately written \path{classify22/check.py} decodes
all $263$ outputs, checks the eigenvector equation with exact
fractions, tests connectivity and pairwise nonisomorphism,
and bijectively matches the independently supplied NetworkX
Graph Atlas cases through order seven. All connected and
regular counts, and the full degree-variety distributions,
agree with the nauty enumeration through order nine and with
the established connected-graph totals.

For order ten, the exhaustive source is the nauty
\texttt{geng} run, not an extrapolation of the C program
(which deliberately stops at nine). Its log records precisely
$11{,}716{,}571$ connected classes, $167$ regular classes and
$599$ nonregular equality classes. The exact independent
per-witness audit yields the displayed subdivision
$538+49+6+6=599$ by degree variety. For $n=2,\ldots,9$
the logs give the other rows exactly. The corresponding
\texttt{SHA256SUMS} files fix the reported program and
output versions. These two enumerations and the exact
row equations constitute the finite certificate.
\end{proof}

One of the order-nine records is the five-degree witness with
column-packed C code \texttt{00d47e9ed}. In vertex order its degrees and
neighbour-degree sums, respectively, are
\[
 (2,3,3,4,4,4,5,5,6),\qquad
 (11,15,15,18,18,18,20,20,21).
\]
Thus $[d(v)^2+2s(v)]/[3d(v)]=13/3$ at each vertex.
The complete edge list is in
\path{certificates/classify22/REPORT.md} and the log; the
certificate independently verifies its connectivity.
This example exists on nine vertices, so the $2k$-vertex
construction of Theorem~\ref{thm:stairs} is not asserted to be
order-minimal. The guide also independently reconstructed the
staircase from its written rule, checking exactly each
$2\leq k\leq40$; these checks corroborate, but do not
replace, the all-$k$ proof.

\section{Prior-art boundaries and open problems}\label{sec:priority}
The dated audit underlying this manuscript is reproduced in
\path{PRIOR-ART.md}. The journal version of
Nikiforov~\cite{Nikiforov2017}, Proposition~20, and his
preprint arXiv:1607.03015 were distinguished; the latter's
printed pages 12--14 contain equations (20)--(21) and
Problem~22. The journal also contains a corrected version
of a \emph{different} least-eigenvalue Proposition~29; that
preprint discrepancy is not relevant to the present result
and is not claimed as a correction to the journal.
Explicit publications of the general integral recognition criterion,
the all-interior-rationals construction, the staircase family, and
the order-ten census were \emph{not found in two independent bounded
audits and the author's search}, all dated 15 November 2026; this
is not proof of publication priority or originality. In particular,
Feng and Yu~\cite[Lemma~2.3]{FengYu2009} established the
$\alpha=1/2$ case, Nikiforov proved regularity for $\alpha>1/2$,
and Alhevaz et al.~\cite[Remark~2.4]{AlhevazEtAl2023} published the
$K_4-e$ expressions whose intersection gives the $\alpha=2/5$
example. Ghebleh et al.~\cite{GheblehEtAl2025} obtain unbounded
degree variety for the distinct, two-main-adjacency-eigenvalue problem;
their construction does not assert equality in (20)--(21) at positive
$\alpha$. The separate access levels and non-exhaustive searches are
documented in \path{PRIOR-ART.md}.
The finite census has not been independently reproduced by
a \emph{third} exhaustive generator at order ten, and no
formal proof-assistant verification is claimed. Floating-point
spectral scans in the background repository are corroboration
only; they do not enter any proof here.

\begin{question}\label{ques:max}
Let $M(n)$ be the maximum number of distinct degrees in a
connected nonregular graph of order $n$ attaining equality
in (20)--(21) at some $0<\alpha<1$ (when such a graph
exists). What is $M(n)$? In particular, what is the
sharp order required for $k$ distinct degrees?
\end{question}
The census gives $M(n)=\lceil n/2\rceil$ for $3\leq n\leq10$.
Theorem~\ref{thm:stairs} proves $M(2k)\geq k$, but does not
prove a universal upper bound or a formula for all $n$.
It also remains open here whether degree partitions with
four degrees can fail to be equitable: none does so in the
complete census through ten vertices, whereas
Proposition~\ref{prop:noneq} shows failure already for
three degrees. Other questions are whether unbounded degree
variety can be realized at each prescribed rational
$\alpha\ne1/3$, and how to classify non-equitable solutions
of \eqref{eq:integer-criterion} without enumeration.
Neither disconnected graphs nor the edge bounds (22)--(23)
of Nikiforov's Problem~22 are classified by the present work.

\subsection*{Acknowledgements and responsibility}
The results were found, developed and checked with AI assistance:
OpenAI ChatGPT research agents contributed proof exploration and code
in separate research threads, while Anthropic's Claude served as
orchestrator and verifier, reading arguments and organizing independent
checks. Aleksei Kudriashov, as author of record for Athanor lab,
takes responsibility for the mathematical claims, citations,
reproducibility statements and final manuscript.

\bibliographystyle{amsplain}
\bibliography{refs}
\end{document}
